\section{Proof of the Fixed-Dimension Minimax Rate}\label{app:minimax}

\begin{proof}[Proof of \Cref{thm:minimax}]
We prove the lower and upper bounds under the same complete-trajectory oracle and the same expected trajectory budget.  For the lower bound, it is enough to construct a two-point submodel contained in $\mathfrak C$ whose CPT gradients are separated by order $\mathsf B^{-1/2}$ while the stopped transcript distributions remain at bounded total variation.  For the upper bound, \Cref{thm:estimator} supplies an unbiased finite-variance output with finite expected trajectory cost.

\paragraph{Step 1: a CPT-realizable one-stock submodel.}
Consider one risky asset and cash, horizon $h=1$, zero transaction cost, and episode-start wealth and reference equal to one.  Take the previous portfolio to be all cash and let the feasibility map be the identity.  Set the risky feature to $e_1\in\R^d$ and the cash feature to zero.  At a query $\theta$, the latent risky weight satisfies
\[
 \mathsf U\sim\operatorname{Beta}(e^{\theta^{[1]}},1),
\]
where $\theta^{[1]}$ is the first coordinate of $\theta$.  After the inertia map \eqref{eq:action-map}, the executed risky weight is $\chi\mathsf U$.

Let the risky return be $R=R_\star\xi$, where $R_\star>0$ is fixed and sufficiently small and
$\xi\sim\operatorname{Bernoulli}(p)$ is independent of $\mathsf U$.  The terminal wealth and reference are
\[
 X_{\rm end}=1+\chi R_\star\mathsf U\xi,
 \qquad
 r_{\rm end}=1+\eta_+\chi R_\star\mathsf U\xi,
\]
so the terminal relative outcome is
\begin{equation}\label{eq:minimax-relative-outcome}
 Y=\upsilon\mathsf U\xi,
 \qquad
 \upsilon=(1-\eta_+)\chi R_\star>0.
\end{equation}
This is an admissible portfolio trajectory model under Definition \ref{def:minimax-oracle}.  The loss-side value is identically zero.

\paragraph{Step 2: local separation of the CPT gradient.}
Let $v=v_+$ and define
\[
 \mathfrak m(t)=\upsilon\,v'(\upsilon t),\qquad 0\le t\le1.
\]
The shifted-power value function gives $\mathfrak m(0)=0$, $\mathfrak m(t)>0$ for $t\in(0,1]$, and bounded $\mathfrak m$.  Since
$\Pp(\mathsf U>t)=1-t^{e^{\theta^{[1]}}}$, a change of variables $y=v(\upsilon t)$ in the gain-side CPT integral gives
\begin{equation}\label{eq:minimax-cpt-submodel}
 J_p(\theta)
 =\int_0^1 \mathfrak m(t)
 w_+^\eps\!\left(p\{1-t^{e^{\theta^{[1]}}}\}\right)dt.
\end{equation}
At the target $\theta_0=0$, differentiation under the integral is valid because, on a compact neighborhood of $\theta_0$,
$t^{e^{\theta^{[1]}}}\abs{\log t}$ has an integrable envelope on $(0,1)$ and $(w_+^\eps)'$ is bounded.  Therefore
\begin{equation}\label{eq:minimax-gradient-function}
 \nabla J_p(0)=\mathfrak g(p)e_1,
 \qquad
 \mathfrak g(p)
 =p\int_0^1\mathfrak m(t)(w_+^\eps)'(p(1-t))(-t\log t)\,dt.
\end{equation}
Let
\[
 d_W=(w_+^\eps)'(0)>0.
\]
Since $\abs{(w_+^\eps)''}\le C_2$, differentiation of \eqref{eq:minimax-gradient-function} with respect to $p$ gives
\begin{align*}
 \mathfrak g'(p)
 &=\int_0^1\mathfrak m(t)(-t\log t)
 \left[(w_+^\eps)'(p(1-t))+p(1-t)(w_+^\eps)''(p(1-t))\right]dt\\
 &\ge (d_W-2C_2p)
 \int_0^1\mathfrak m(t)(-t\log t)\,dt.
\end{align*}
The integral is strictly positive.  Choose $p_{\max}\in(0,1/2]$ small enough that $2C_2p_{\max}\le d_W/2$; if $C_2=0$, any choice in $(0,1/2]$ is admissible.  Let
\[
 p_*=p_{\max}/2,
 \qquad
 \mathcal I=[p_*/2,3p_*/2],
 \qquad
 c_*=
 \frac{d_W}{2}
 \int_0^1\mathfrak m(t)(-t\log t)\,dt.
\]
Then
\begin{equation}\label{eq:minimax-gradient-separation}
 \mathfrak g'(p)\ge c_*>0
 \qquad\text{for every }p\in\mathcal I.
\end{equation}

\paragraph{Step 3: information in an adaptive stopped transcript.}
For $\mathsf B$ large enough, take
\[
 p_\pm=p_*\pm\delta,
 \qquad
 \delta=\delta_*/\sqrt{\mathsf B},
\]
where $\delta_*>0$ is chosen so that $p_\pm\in\mathcal I$.  For Bernoulli laws in the fixed interior interval $\mathcal I$, the elementary bound
$\operatorname{kl}(p\Vert q)\le (p-q)^2/\{q(1-q)\}$ gives
\begin{equation}\label{eq:bernoulli-kl-bound}
 \operatorname{kl}(p_+\Vert p_-)
 \le C_I\delta^2
\end{equation}
for a finite constant $C_I$ depending only on $\mathcal I$.

Consider an arbitrary admissible estimator in \eqref{eq:minimax-risk}.  Its query point at each oracle call may depend on the preceding transcript.  Conditional on any such query, the latent Dirichlet action law is identical under $p_+$ and $p_-$; only the Bernoulli return kernel differs.  All algorithmic randomization can be included in the transcript and is likewise common to the two models.  Let $\Pp_+^{\rm tr}$ and $\Pp_-^{\rm tr}$ denote the resulting stopped transcript laws.  The data-dependent call count $T_{\rm call}$ is a stopping time for the oracle transcript filtration.  The chain rule for the stopped likelihood ratio gives
\begin{equation}\label{eq:stopped-transcript-kl}
 \operatorname{KL}(\Pp_+^{\rm tr}\Vert\Pp_-^{\rm tr})
 =\E_+T_{\rm call}\,\operatorname{kl}(p_+\Vert p_-)
 \le \mathsf B C_I\delta^2=C_I\delta_*^2.
\end{equation}
The equality can be obtained by summing the conditional log-likelihood increments over
$\{T_{\rm call}\ge j\}$; these events are measurable before the $j$th fresh return is observed.  The increments are uniformly integrable because $p_\pm$ stay in the compact interior interval $\mathcal I$ and $\E_+T_{\rm call}<\infty$.

Choose $\delta_*$ small enough that $C_I\delta_*^2\le1/8$.  Pinsker's inequality then gives
\begin{equation}\label{eq:minimax-tv}
 \operatorname{TV}(\Pp_+^{\rm tr},\Pp_-^{\rm tr})\le\frac14.
\end{equation}

\paragraph{Step 4: reduction from estimation to testing.}
By \eqref{eq:minimax-gradient-separation} and the mean-value theorem, the two gradient targets satisfy
\begin{equation}\label{eq:minimax-target-distance}
 \Delta_{\mathsf B}^g
 :=\norm{\mathfrak g(p_+)e_1-\mathfrak g(p_-)e_1}_2
 \ge2c_*\delta.
\end{equation}
Given any estimator $\widehat g$, define the test that chooses the closer of the two targets.  On a testing error, the estimator must be at least $\Delta_{\mathsf B}^g/2$ from the true target.  The sum of the two testing error probabilities is at least one minus the total variation distance.  Therefore
\begin{align*}
 \max_{s\in\{+,-\}}
 \E_s\norm{\widehat g-\mathfrak g(p_s)e_1}_2^2
 &\ge \frac{(\Delta_{\mathsf B}^g)^2}{8}
 \left\{1-\operatorname{TV}(\Pp_+^{\rm tr},\Pp_-^{\rm tr})\right\}\\
 &\ge \frac{3c_*^2\delta_*^2}{8\mathsf B}.
\end{align*}
The two models belong to $\mathfrak C$, so taking the infimum over all admissible procedures gives the lower bound in \eqref{eq:minimax-rate} with
$c_{\rm lb}=3c_*^2\delta_*^2/8$.

\paragraph{Step 5: cost-matched upper bound.}
Fix any admissible estimator design parameters $(n,p_{\rm act},b)$ from \Cref{thm:estimator}.  The common model-class constants in Definition \ref{def:minimax-oracle} make
\[
 \sup_{P\in\mathfrak C}\tr\Cov_P(Z_n^{\db})\le C_{\rm var}(n)<\infty,
 \qquad
 C_{\rm cost}=n(1+p_{\rm act} \Upsilon_b)<\infty
\]
uniform over $P\in\mathfrak C$.  For
\[
 M=\left\lfloor\frac{\mathsf B}{C_{\rm cost}}\right\rfloor,
\]
run $M$ independent, untruncated debiased outputs $Z_n^{\db}$ at $\theta_0$ and average them.  Their total expected number of complete trajectories is $MC_{\rm cost}\le \mathsf B$, so this is an admissible procedure.  For $\mathsf B\ge2C_{\rm cost}$,
$M\ge \mathsf B/(2C_{\rm cost})$.  Hence
\[
 \sup_{P\in\mathfrak C}
 \E_P\norm{\widehat g_{M,n}^{\db}-g_P}_2^2
 \le\frac{C_{\rm var}(n)}{M}
 \le\frac{2C_{\rm var}(n)C_{\rm cost}}{\mathsf B}.
\]
Thus the upper bound in \eqref{eq:minimax-rate} holds with
$C_{\rm ub}=2C_{\rm var}(n)C_{\rm cost}$ and a sufficiently large $\mathsf B_0$.  The lower and upper bounds use the same target, trajectory oracle, loss, and expected raw-trajectory budget, completing the proof.
\end{proof}
